\scanpage{196}
\mspage{L 19}
% archive sheet number: 195

que en l'occurence) ne s'apparente pas au pur travail du marteau sur l'enclume ou sur le burin\deleted{\unsure{er|e}} - le mépris des processus créateurs \inserted{\handwritten{les}} plus délicats (et \inserted{\handwritten{parfois}} de moindre apparence) , \deleted{et à la limite} de tout ce qui est \deleted{idée, vision,} inspiration, rêve , vision \deleted{ou \deleted{même} simplement idée, même si clairement conçue et formulée soit-elle} \inserted{\handwritten{si puissants et fertiles soient-ils}} \inserted{\handwritten{et même}} (à la limite) de toute \mbox{i\,d\,é\,e} , si clairement conçue et formulée soit-elle; de tout ce qui n'est écrit et \mbox{p\deleted{\unsure{u}}u\,b\,l\,i\,é} \unsure{noir|noit} sur blanc, sous forme d'énoncés purs et durs, \deleted{et avec démontrés,} \unsure{répertoriables|répertorizables|répertoriéables} et répertoriés, mûrs pour les ``banques de \uline{données''} s'engouffrant dans \deleted{\unsure{d'}}\inserted{les} inépuisables mémoires de nos mégaordinateurs.

Il y a eu \deleted{\illegible{1}} (pour reprendre \deleted{les} une expression de C.L. Siegel ($^{\ast}$)) un extraordinaire ``\uline{aplatissement}'' ou ``\uline{rétrécissement}'' dans la pensée mathématique($^{\ast}$), dépouillée d'une dimension essentielle, \deleted{d'un ``versant''} \inserted{\handwritten{de \deleted{tout} son}} \deleted{de son} ``versant\deleted{\unsure{''}} \inserted{\handwritten{d'ombre''}}, \deleted{\illegible{1}} \inserted{\handwritten{du versant}} \deleted{\illegible{1}} ``féminin''. \deleted{Traditionnellement,} \inserted{\handwritten{Il est vrai que de tradition ancestrale,}} ce versant-là \deleted{restait} \inserted{du travail de découverte} \deleted{scientifique} \inserted{restait} dans une large mesure occulté, \deleted{\unsure{dans la on n'en|dans les on n'en}} personne (autant dire) n'en \uline{parlait} jamais - mais le contact avec les sources profondes du rêve, qui alimentent \deleted{les grands desseins et} les grandes visions et les grands desseins, \deleted{\unsure{n'a jamais nté|n'a jamais été}} \inserted{n'avait} jamais \inserted{\handwritten{encore}} (à ma connaissance) était perdu. Il semblerait \deleted{que} \inserted{qu'à présent} nous sommes déjà entrés dans une époque \deleted{où ce contact est sur le point d'être} \inserted{de \mbox{d\,e\,s\,s\,é\,c\,h\,e\,m\,e\,n\,t},} \inserted{où} cette source est, non point tarie certes, mais où l'accès à elle est condamné\deleted{\unsure{e}}, par le verdict sans appel \deleted{de la dérision ed} du mépris \deleted{généraux} \inserted{général} et par les représailles de la dérision. Nous approchons du moment, \unsure{semble-tèil|semble-t-il}, où \inserted{\unsure{sera éradiqué|sesa éradiqué}} non seulement le \uline{souvenir} \inserted{\handwritten{en chacun}} \deleted{des stades et étapes} \inserted{de \deleted{ces} tout travail} \deleted{\unsure{à notre|d'une}} proche de la source, du travail \handwritten{à} \deleted{\unsure{lien à|lien a}} dominante ``féminine'', \deleted{dans le travail mathématique (stades} (ridiculisés par des épithètes comme ``vaseux'', ``mou\deleted{\unsure{s}}'', ``inconsistant\deleted{\unsure{s}}'' - \deleted{ou} au bout opposé, \deleted{``trivial'', ``évident'',} \inserted{``trivialités'', ``évidences'',} ``bombinage'' \deleted{de} ...), mais \inserted{\handwritten{où sera extirpé}} \deleted{des stades} \inserted{ce travail} \deleted{étapes elles-mêmes, celles qui où naissent les} \inserted{de gestation même: celui} où sont conçues, s'élaborent et naissent les notions et les visions nouvelles. \deleted{\unsure{Le} \handwritten{\unsure{que}}} \inserted{Ce sera} \inserted{aussi} où l'exercice de notre art sera réduit à \deleted{\unsure{un|de}}\handwritten{\unsure{de|un}} \inserted{\unsure{aride|ar}\handwritten{\unsure{de|le}}} exercice de ``poids et haltères''\inserted{cérébral} - ce\handwritten{lle} d'une hypertrophie \inserted{fiévreuse et} stérile, prenant la suite de plus de trois siècles \deleted{\unsure{de|ce}}\inserted{\handwritten{\unsure{d'i|dit|d'}}} renouvellement créateur...

Mais à nouveau je \deleted{anticipe ou} digresse, \deleted{\unsure{ou}} j'anticipe sur ce que la réflexion m'a enseigné, \deleted{dans la en continu} \inserted{\handwritten{en partant}} \deleted{dans la direction} d'un double propos clairement présent en moi \deleted{en commençant mon ``introduction'', en juin 1983 d'abord et en la reprenant en février l'année dernier (l'année d'après):} \inserted{\handwritten{dès avant même les débuts de \deleted{\illegible{1}} celui-ci:}} le propos d'une ``déclaration d'intentions'', et (intimement lié à celui-ci) celui de m'exprimer au sujet de \deleted{\unsure{b}}\unsure{la} nature du travail créateur. Il y avait pourtant un troisième propos encore, moins clairement présent sûrement au niveau conscient, mais correspondant à un besoin plus profond
